\documentclass[11pt,a4paper]{article}

% ----------------------------------------------------------------------
% Packages
% ----------------------------------------------------------------------
\usepackage[utf8]{inputenc}
\usepackage[T1]{fontenc}
\usepackage[a4paper,margin=2.5cm]{geometry}
\usepackage{graphicx}
\usepackage{xcolor}
\usepackage{booktabs}
\usepackage{enumitem}
\usepackage{titlesec}
\usepackage{fancyhdr}
\usepackage{parskip}
\usepackage[hidelinks]{hyperref}

% ----------------------------------------------------------------------
% Formatting
% ----------------------------------------------------------------------
\definecolor{navy}{RGB}{46,64,87}
\titleformat{\section}{\Large\bfseries\color{navy}}{\thesection}{1em}{}
\titleformat{\subsection}{\large\bfseries\color{navy}}{\thesubsection}{1em}{}

\pagestyle{fancy}
\fancyhf{}
\fancyhead[L]{\small SOC Home Lab with SIEM}
\fancyhead[R]{\small M. A. Khalid}
\fancyfoot[C]{\thepage}
\renewcommand{\headrulewidth}{0.4pt}

\setlength{\parskip}{0.6em}

% ----------------------------------------------------------------------
% Title page content
% ----------------------------------------------------------------------
\begin{document}

\begin{titlepage}
    \centering
    \vspace*{2cm}

    {\Huge\bfseries\color{navy} Building a Reproducible Security\\[0.3cm]
    Operations Centre Home Lab with an\\[0.3cm]
    Open-Source SIEM\par}

    \vspace{1.5cm}
    {\Large A Practical Study in Security Monitoring,\\ Endpoint Telemetry, and Alert Triage\par}

    \vspace{2.5cm}
    {\large\textbf{Muhammad Ahmed Khalid}\par}
    \vspace{0.3cm}
    {\large Personal Portfolio Project\par}

    \vspace{2cm}
    \begin{minipage}{0.85\textwidth}
        \centering\itshape\small
        This report documents a self-initiated technical project. It is not
        submitted for academic credit. It accompanies the project repository
        and is intended to demonstrate practical security-monitoring skills.
    \end{minipage}

    \vfill
    {\large \today\par}
\end{titlepage}

% ----------------------------------------------------------------------
% Abstract
% ----------------------------------------------------------------------
\section*{Abstract}
\addcontentsline{toc}{section}{Abstract}

Security Operations Centre (SOC) analysts depend on Security Information and
Event Management (SIEM) platforms to detect and investigate threats, yet
graduates entering the field rarely have hands-on experience with these tools.
This project addresses that gap by designing, building, and documenting a
fully functional SOC home lab using the open-source SIEM platform Wazuh.

The lab consists of a SIEM server and two monitored endpoints --- a Windows
host and a Linux host --- deployed on an isolated virtual network. Enhanced
endpoint telemetry is provided by Sysmon on Windows and the Linux Audit
framework on Linux. The environment is used to generate security events,
observe the resulting alerts, and demonstrate a realistic analyst triage
workflow. A core aim of the project is reproducibility: every step is
documented so a third party can rebuild the lab from the accompanying
repository.

\vspace{1em}
\noindent\textit{Keywords:} SIEM, Security Operations Centre, Wazuh, threat
detection, endpoint telemetry, log analysis, home lab.

% ----------------------------------------------------------------------
% Table of contents
% ----------------------------------------------------------------------
\newpage
\tableofcontents
\newpage

% ----------------------------------------------------------------------
% Sections
% ----------------------------------------------------------------------
\section{Introduction}
\label{sec:introduction}

Modern organisations generate enormous volumes of security-relevant data:
authentication events, process executions, network connections, file changes,
and system errors. Making sense of this data --- and acting on it quickly
enough to matter --- is the work of a Security Operations Centre (SOC). At the
centre of every SOC sits a Security Information and Event Management (SIEM)
platform, which collects logs from across an estate, correlates them, and
raises alerts an analyst can investigate.

For someone entering the cybersecurity profession, this creates a practical
difficulty. Almost every junior SOC and security-analyst role expects
familiarity with a SIEM, yet undergraduate study tends to teach the theory of
detection and incident response without sustained hands-on exposure to the
tools themselves. A graduate can describe what a SIEM does without ever having
operated one.

This project was undertaken to close that gap directly. It designs, builds,
and documents a working SOC home lab using only free and open-source software,
producing both a functional environment and a reproducible record of how it
was created.

\subsection{Aim and Objectives}
\label{sec:aim}

The \textbf{aim} of this project is to design, deploy, and operate a functional
SOC home lab built on an open-source SIEM, capable of collecting and analysing
security telemetry from multiple endpoints and supporting a realistic analyst
triage workflow.

The aim is met through the following \textbf{objectives}:

\begin{enumerate}[label=\textbf{O\arabic*.},leftmargin=*]
    \item Deploy an open-source SIEM server on an isolated virtual network.
    \item Configure two monitored endpoints --- a Windows host and a Linux
          host --- each reporting to the SIEM.
    \item Enable enhanced host telemetry beyond default operating-system logs.
    \item Verify centralised log collection and establish a baseline of normal
          activity.
    \item Simulate a set of security events and observe the resulting alerts.
    \item Triage the alerts, classifying them and documenting the reasoning.
    \item Produce a fully reproducible record of the build so the lab can be
          rebuilt by a third party.
\end{enumerate}

\subsection{Research Questions}
\label{sec:rq}

The project is framed around three research questions, revisited in the
Discussion and Conclusions:

\begin{itemize}[leftmargin=*]
    \item \textbf{RQ1.} Can a credible SOC monitoring environment, capable of
          collecting endpoint telemetry and generating meaningful alerts, be
          built using only free and open-source tools?
    \item \textbf{RQ2.} What additional configuration is required beyond a
          default SIEM and operating-system installation to achieve
          detection-grade visibility?
    \item \textbf{RQ3.} To what extent does a single-machine virtualised lab
          realistically represent the work of a professional SOC, and where
          are the limits of that representation?
\end{itemize}

\subsection{Contributions}
\label{sec:contributions}

This project contributes:

\begin{itemize}[leftmargin=*]
    \item A complete, working SOC home-lab build using Wazuh, deployed on an
          isolated network with Windows and Linux endpoints.
    \item A fully documented, step-by-step record enabling independent
          reproduction of the environment.
    \item A set of attack simulations paired with the alerts they produce and
          an analyst-style triage of each.
    \item An honest assessment of how closely such a lab represents real SOC
          work, including its limitations.
\end{itemize}

\subsection{Report Structure}
\label{sec:structure}

The remainder of this report is organised as follows.
\textbf{Section~\ref{sec:litreview}} reviews the background literature on
SOCs, SIEM technology, logging, and security home labs, and positions this
work against related projects.
\textbf{Section~\ref{sec:methodology}} sets out the general methodology and
design decisions.
\textbf{Section~\ref{sec:implementation}} documents the specific
implementation of the lab.
\textbf{Section~\ref{sec:evaluation}} evaluates whether the lab meets its
objectives.
\textbf{Section~\ref{sec:discussion}} discusses challenges and revisits the
research questions, and \textbf{Section~\ref{sec:conclusions}} concludes and
identifies future work.

\section{Literature Review and Background Knowledge}
\label{sec:litreview}

\subsection{Overview}

This section establishes the background needed to understand the project. It
begins with the function of a Security Operations Centre and the role of the
analyst, then examines the technology that underpins it: Security Information
and Event Management (SIEM) platforms, log management, and endpoint telemetry.
It compares open-source and commercial SIEM options, outlines the main
approaches to threat detection and the MITRE ATT\&CK framework, and considers
the use of virtualised home labs for developing practical skills. The section
closes with a review of related work and a statement of how this project
differs from it.

\subsection{Security Operations Centres}

A Security Operations Centre (SOC) is the organisational function responsible
for continuously monitoring an environment, detecting malicious or anomalous
activity, and coordinating the response to security incidents. Rather than a
single technology, a SOC is best understood as the combination of people,
process, and tooling working together to maintain an organisation's security
posture.

SOC analysts are commonly organised into tiers. Tier~1 analysts perform the
initial triage of incoming alerts, deciding quickly whether an alert
represents genuine malicious activity (a true positive) or harmless behaviour
that merely resembles it (a false positive). Alerts that cannot be dismissed
are escalated to Tier~2 analysts for deeper investigation, and the most
serious or complex cases reach Tier~3 staff who carry out incident response
and proactive threat hunting. This triage-first model means that the largest
share of day-to-day SOC work consists of rapidly reviewing, classifying, and
documenting alerts.

The structured handling of incidents that a SOC supports is described in
NIST's incident-handling guidance, which frames incident response as a
lifecycle of preparation; detection and analysis; containment, eradication and
recovery; and post-incident activity~\cite{nist80061}. Detection and analysis
--- the stage at which raw data becomes an actionable alert --- is the stage
this project is most concerned with, and it depends almost entirely on the
SOC's central tooling: the SIEM.

\subsection{Security Information and Event Management}

A SIEM platform is the technological core of a SOC. The term combines two
earlier categories: Security Information Management (SIM), concerned with the
long-term collection, storage, and analysis of log data, and Security Event
Management (SEM), concerned with the real-time monitoring and correlation of
events. A modern SIEM unifies both.

In practice, a SIEM performs several core functions. It \textit{collects} log
and event data from many different sources across an estate. It
\textit{normalises} that data, parsing records written in many different
formats into a consistent structure so they can be compared and searched. It
\textit{correlates} events, applying rules that recognise patterns spanning
multiple log sources --- for example, linking a series of failed logins
followed by a success into a single suspicious sequence. It \textit{alerts}
analysts when activity matches a rule of concern, and it provides
\textit{dashboards and search} so that analysts can investigate. Finally, it
\textit{retains} data for the periods required by forensic investigation and
regulatory compliance.

SIEM platforms have evolved considerably. Early systems were primarily log
aggregators; later generations added correlation and alerting; current
platforms increasingly incorporate user and entity behaviour analytics,
threat-intelligence enrichment, and automated response. Throughout this
evolution, the underlying premise has remained constant: an organisation
cannot detect what it does not collect and analyse, and the SIEM is the
instrument through which that collection and analysis happens.

\subsection{Log Management and Endpoint Telemetry}

Because detection ultimately rests on log data, the quality of that data is
decisive. NIST's guidance on computer security log management emphasises that
logs are essential not only for detecting attacks but also for understanding
them after the fact, and that effective log management requires deliberate
planning of what to collect and how to handle it~\cite{nist80092}. National
guidance makes a similar point: logging must be designed with security
outcomes in mind rather than left to operating-system defaults~\cite{ncsc}.

This distinction --- between default logging and security-grade logging --- is
central to the present project. On Windows, the standard Event Log captures
useful records such as logon events and service changes, but it does not by
itself provide the fine-grained visibility that behavioural detection
requires. \textbf{Sysmon} (System Monitor), a free tool from Microsoft's
Sysinternals suite, addresses this by logging detailed information about
process creation, network connections, image loads, file-creation timestamps,
and related activity, written to a dedicated event channel~\cite{sysmon}. This
additional telemetry is what allows a SIEM to reason about \textit{behaviour}
rather than merely about discrete system messages.

Linux provides an equivalent capability through the Linux Audit framework and
its userspace daemon, \texttt{auditd}. By defining audit rules, an
administrator can record security-relevant system activity --- notably the
execution of commands through the \texttt{execve} system call --- at the
kernel level. As with Sysmon on Windows, this elevates the endpoint from
producing general operational logs to producing detection-grade telemetry.

The recurring principle across both platforms is that endpoint visibility is
something that must be deliberately engineered. A SIEM connected only to
default logs will see far less than one fed by purpose-configured endpoint
telemetry.

\subsection{Open-Source SIEM Platforms}

SIEM capability is available both as commercial products and as open-source
software. Commercial platforms such as Splunk, Microsoft Sentinel, and IBM
QRadar are powerful and widely deployed, but they are typically licensed by
data volume, which can make them costly --- and largely inaccessible for
independent learning and experimentation.

Several capable open-source alternatives exist, including the Elastic Stack
(often deployed as Elastic Security), Security Onion, Graylog, and Wazuh. This
project uses \textbf{Wazuh}. Wazuh originated as a fork of
\textbf{OSSEC}, a long-established open-source host-based intrusion detection
system~\cite{ossec}, and has since grown into a full SIEM and extended
detection and response (XDR) platform~\cite{wazuhdocs}. A Wazuh deployment
consists of a central \textit{manager} that analyses incoming data and applies
detection rules, an \textit{indexer} that stores and searches events, a
\textit{dashboard} that presents them, and lightweight \textit{agents}
installed on monitored endpoints. Beyond log analysis and alerting, Wazuh
provides file-integrity monitoring, configuration assessment, vulnerability
detection, and the mapping of alerts to the MITRE ATT\&CK framework.

For the purpose of building practical skills, an open-source platform offers
clear advantages. It can be deployed at no cost, it can be examined and
reconfigured freely, and --- importantly --- the concepts it teaches, such as
agent-based collection, log parsing, rule-based correlation, and alert triage,
transfer directly to commercial platforms. The specific product differs; the
underlying discipline does not.

\subsection{Threat Detection Approaches and the MITRE ATT\&CK Framework}

SIEM platforms detect threats using two broad approaches.
\textbf{Signature- or rule-based} detection matches activity against known
patterns of malicious behaviour; it is precise and explainable, but it can
only catch what has been defined in advance. \textbf{Anomaly-based} detection
flags deviations from an established baseline of normal behaviour; it can in
principle catch novel activity, but it is prone to raising alerts on benign
deviations. NIST's guidance on intrusion detection and prevention systems
describes both approaches and notes that practical deployments generally
combine them, since each compensates for the other's
weakness~\cite{nist80094}.

A persistent operational problem for SOCs is \textit{alert fatigue}: when
detection produces a high volume of low-value or false-positive alerts,
analysts become desensitised, and genuine threats risk being overlooked among
the noise. This makes the \textit{tuning} of detection --- raising true
positives while suppressing false positives --- as important as detection
itself.

To make detection systematic, the security community has increasingly adopted
the \textbf{MITRE ATT\&CK} framework: a structured, continually updated
knowledge base of adversary tactics and techniques drawn from real-world
observations of attacks~\cite{mitreattack}. ATT\&CK gives defenders a common
vocabulary and a way to measure detection \textit{coverage} --- mapping which
adversary techniques an environment can and cannot detect. While the detailed
application of ATT\&CK belongs to the planned follow-on project on detection
engineering, it is introduced here because it is the framework against which
modern detection capability is now routinely assessed.

\subsection{Virtualised Security Labs and the Skills Gap}

The cybersecurity profession is widely recognised to face a shortage of
skilled practitioners, and industry reporting consistently identifies a
mismatch between the volume of security-relevant activity organisations must
handle and the workforce available to handle it~\cite{verizondbir}. For new
entrants specifically, the gap is often less about theoretical knowledge than
about practical, tool-level experience.

Virtualised home labs are a well-established response to this problem. Using a
hypervisor such as VirtualBox or VMware, a learner can build a self-contained
environment of multiple virtual machines on a single physical computer, at no
cost and in complete isolation from any production network. Such a lab can be
broken, rebuilt, and experimented on freely. For defensive security in
particular, a home lab allows a learner to stand up a SIEM, generate security
events safely, and practise the monitoring and triage workflow that a SOC role
demands --- experience that is otherwise difficult to obtain before
employment. Frameworks such as the CIS Critical Security Controls reinforce
the value of this hands-on focus by expressing good security practice in terms
of concrete, implementable actions rather than abstract principles~\cite{ciscontrols}.

\subsection{Related Works}

Building security home labs is an established practice, and several reference
projects exist. The most prominent is \textbf{DetectionLab} by Chris
Long~\cite{detectionlab}, an automated project that provisions a pre-built
Windows domain instrumented with logging and monitoring tooling. DetectionLab
is comprehensive and highly automated; its strength --- one-command
provisioning --- is also a limitation for a learner, as the automation
abstracts away the very build steps from which much of the understanding
comes. A large body of community material, including vendor tutorials and
blog write-ups, similarly documents Wazuh and Elastic-based lab builds, though
such material varies in depth and is rarely presented as a structured,
reproducible whole.

This project is positioned deliberately against that backdrop, and differs in
three respects. First, \textbf{the build is performed and documented manually
and in full}: reproducibility is treated as a primary deliverable, with each
step recorded so that the understanding --- not only the result --- is
transferable. Second, \textbf{the evaluation is explicitly honest about
realism}: rather than presenting a home lab as equivalent to a production SOC,
the report sets out clearly where a single-machine virtualised environment
does and does not represent real operations. Third, \textbf{the project is
designed as the first stage of a two-part progression}: it establishes the
monitoring foundation on which a planned second project, focused on detection
engineering with MITRE ATT\&CK and Atomic Red Team~\cite{atomicredteam}, will
build. The contribution is therefore not a novel tool, but a clear,
honest, and reproducible account of how a credible SOC monitoring environment
is constructed and operated using open-source software.

% NOTE TO AUTHOR: The Related Works subsection currently cites only sources
% that can be verified. To strengthen it for an academic standard, add 2-3
% peer-reviewed papers (e.g. studies evaluating Wazuh or open-source SIEM
% detection performance) found through your own literature search. Do not add
% any reference you have not personally read and verified.

\section{Methodology}
\label{sec:methodology}

\subsection{Overview}

This section sets out the general approach taken to the project. It covers the design principles that shaped every subsequent decision, the build sequence philosophy, the simulation methodology used to validate the lab, and the treatment of documentation as a deliverable in its own right. Specific commands, configurations, and technical detail are deferred to Section~\ref{sec:implementation}.

\subsection{Guiding Design Principles}

Four principles were fixed at the start of the project and applied consistently to every subsequent decision.

\textbf{Isolation.} The lab must never be able to reach or be reached by any production network or the wider internet in a way that permits the escape of simulated attack activity. This is fundamental for any environment in which security incidents are deliberately generated. It was addressed by deploying every virtual machine onto a dedicated VirtualBox NAT Network (\texttt{SOC-Lab}, subnet 10.0.10.0/24), with dashboard access from the host provided only through an explicit port-forwarding rule. No bridged adapters were used.

\textbf{Open-source, reproducible tooling.} All security technology used in the lab is open source and free to obtain. This is not merely a cost decision; it also means that a third party can reproduce the environment from scratch without any licensing barrier, and that the underlying operation of each component can be inspected. Wazuh was selected as the SIEM, Sysmon as the Windows telemetry source, and the Linux Audit framework (\texttt{auditd}) as the Linux equivalent.

\textbf{Reproducibility as a first-class deliverable.} Every step of the build was designed to be documented at the time it was performed, with the intention that the documentation would enable a third party to rebuild the environment from a clean state. Documentation was therefore not a task performed at the end of the project, but an ongoing output produced alongside every configuration change.

\textbf{Honest evaluation.} A home lab cannot fully reproduce the scale, noise, or operational realities of a production SOC. Rather than obscure this, the project committed at the outset to evaluating its own limitations openly, on the basis that transparency about what a training environment can and cannot represent is more useful than an inflated claim of realism.

\subsection{Build Sequence}

The build followed an iterative pattern of \textit{build, test, document}, applied at each of seven milestones:

\begin{enumerate}[label=\textbf{M\arabic*.},leftmargin=*]
    \item Environment and network setup.
    \item Wazuh server deployment and dashboard access.
    \item Windows endpoint build, agent enrollment, and Sysmon configuration.
    \item Linux endpoint build, agent enrollment, and \texttt{auditd} configuration.
    \item Attack simulations across both endpoints.
    \item Analyst triage and incident report writing.
    \item Formal report authoring.
\end{enumerate}

Each milestone was completed before the next was started, and each produced its own documentation artefact. A component was considered complete only when the corresponding documentation was written to a standard that would allow reconstruction of that step by someone reading only the document.

\subsection{Simulation Methodology}

The attack simulations in milestone M5 were selected to satisfy three criteria simultaneously. First, they had to be safe: no actual malware, no attempts to escape the isolated network, no destructive actions on the host. Second, they had to be realistic: each simulation was chosen to represent a well documented adversary technique, ideally mappable to a MITRE ATT\&CK identifier, so that the resulting detections could be reasoned about in the same language a real SOC would use. Third, they had to span multiple detection classes so that the resulting alert coverage was not concentrated in a single Wazuh rule family.

Four simulations were executed in total, targeting:

\begin{itemize}[leftmargin=*]
    \item Windows local account creation and privilege escalation (T1136.001, T1078.003).
    \item Linux local account creation and privilege escalation via \texttt{sudo}.
    \item Repeated failed \texttt{sudo} authentication attempts (T1110, credential access via brute force).
    \item Windows hosts file modification, targeting File Integrity Monitoring capabilities.
\end{itemize}

For each simulation, evidence was captured in the form of dashboard screenshots showing the resulting Wazuh alerts, and each was reviewed as an analyst would review a real alert cluster in production.

\subsection{Analyst Workflow}

The project treated the analyst review of resulting alerts as a discipline in its own right, not merely a validation step. For the strongest of the four simulations, a full incident report was written in a professional format including MITRE ATT\&CK mapping, timeline reconstruction, technical analysis, impact assessment, response actions, and recommendations. This was intentional: writing up a simulated incident in the same form a real one would take is what turns a technical exercise into meaningful analyst practice.

\subsection{Documentation Approach}

Documentation for the project takes three forms. First, per-milestone Markdown documents in the project's \texttt{docs/} folder describe the build of each component, versioned in Git alongside the code. Second, dashboard screenshots for every simulation are stored in the \texttt{screenshots/} folder and referenced from the incident report. Third, this formal report (produced in \LaTeX) captures the project as a whole, including methodology, evaluation, and honest reflection on limitations. All three are published to a public GitHub repository so that the project is reviewable in full by any third party.

\section{Implementation}
\label{sec:implementation}

\subsection{Overview}

This section walks through the actual build of the lab: the network, the SIEM
manager, the two endpoints, the endpoint telemetry enhancements, and the
verification that logs flow end to end. It corresponds to the full
step-by-step documentation stored in the project repository under
\texttt{docs/}, distilled here into narrative form. The precise commands,
configuration snippets, and screenshots referenced in this section are held in
those source documents so this report can remain readable while still being
independently reproducible.

\subsection{Network and Virtualisation}
\label{sec:impl-network}

Oracle VirtualBox 7.2.8 was used to host all virtual machines. A dedicated
\textit{NAT network} named \texttt{SOC-Lab} was created on the address range
\texttt{10.0.10.0/24} as described in Section~\ref{sec:env-design}. The NAT
network mode was chosen deliberately over the more permissive Bridged mode:
VMs on a NAT network can reach each other but cannot be reached from, or
directly reach, the host network without deliberate port forwarding. This
provides the isolation guarantee that a security lab must offer by default.

Three virtual machines were provisioned onto this network:

\begin{itemize}[leftmargin=*]
    \item The Wazuh server VM at \texttt{10.0.10.3}.
    \item The Windows endpoint VM at \texttt{10.0.10.4}.
    \item The Linux endpoint VM at \texttt{10.0.10.5}.
\end{itemize}

Fixed addressing was preferred over DHCP for a lab of this size, so that
sources in the Wazuh dashboard would always resolve to the same host without
needing to cross-reference DHCP leases.

\subsection{Wazuh SIEM Deployment}
\label{sec:impl-wazuh}

The Wazuh server was deployed from the official
\texttt{wazuh-4.14.5.ova} appliance provided by Wazuh Inc. The appliance is
a pre-built all-in-one deployment containing the three Wazuh components
required to operate the platform:

\begin{itemize}[leftmargin=*]
    \item \textbf{wazuh-manager} --- the analysis and rule engine that receives
          agent telemetry, applies detection rules, and generates alerts.
    \item \textbf{wazuh-indexer} --- the storage and search layer, built on
          OpenSearch, that persists alerts and events so they can be queried
          and displayed.
    \item \textbf{wazuh-dashboard} --- the browser-based user interface, built
          on OpenSearch Dashboards, that presents alerts and enables analyst
          workflows.
\end{itemize}

Using the appliance rather than performing a from-source install was a
pragmatic choice: it reduced the initial deployment to VM import and network
attachment, which allowed the project to move on to the more instructive
work of endpoint configuration and detection engineering. The trade-off is
that the appliance is opaque compared to a from-source deployment; a
production build would follow Wazuh's documented distributed installation
procedure so component sizing could be tuned per role.

After boot, the appliance's default credentials were rotated, its network
interface was confirmed on \texttt{10.0.10.3}, and the dashboard was
reached at \texttt{https://10.0.10.3} from the host machine. The dashboard's
first-run wizard was used to generate agent deployment commands for the two
endpoints.

\subsection{Windows Endpoint}
\label{sec:impl-windows}

The Windows endpoint was provisioned from Microsoft's \texttt{WinDev2407Eval}
appliance (Windows 11), a free evaluation image intended for developer
testing. This choice avoided the licensing and installation overhead of a
manual Windows install while remaining fully representative of a modern
Windows endpoint.

Three configuration steps were performed after the VM was joined to the
\texttt{SOC-Lab} network:

First, the \textbf{Wazuh agent} was installed using the PowerShell one-liner
generated by the manager's Deploy new agent wizard. The wizard tailored the
command with the manager IP (\texttt{10.0.10.3}) and the assigned agent name
(\texttt{Windows-Endpoint}). The agent service was started with
\texttt{NET START WazuhSvc} and appeared as \texttt{Active} in the
manager's agent list within seconds.

Second, \textbf{Sysmon} was installed with the SwiftOnSecurity baseline
configuration. Sysmon writes to a dedicated Windows event channel,
\texttt{Microsoft-Windows-Sysmon/Operational}, which is separate from the
Windows Security event log and requires explicit inclusion in the Wazuh
agent's configuration to be forwarded. The following block was added to the
agent's \texttt{ossec.conf} inside the \texttt{<ossec\_config>} element:

\begin{quote}\small
\texttt{<localfile>}\\
\texttt{~~<location>Microsoft-Windows-Sysmon/Operational</location>}\\
\texttt{~~<log\_format>eventchannel</log\_format>}\\
\texttt{</localfile>}
\end{quote}

The agent service was restarted, and within one minute the manager began
receiving events under the \texttt{sysmon}, \texttt{sysmon\_event1}, and
\texttt{sysmon\_event\_3} rule groups. This confirmed that the SwiftOnSecurity
configuration was producing telemetry and that the Wazuh manager's Sysmon
decoder was correctly parsing the incoming events.

Third, a \textbf{baseline of normal activity} was observed for approximately
twenty-four hours. This served two purposes: it validated that the agent
would continue reporting stably across VM restarts, and it established a
reference of what "clean" telemetry from the endpoint looked like, so that
subsequent simulation-generated alerts could be distinguished from routine
noise.

Full step-by-step commands, screenshots, and verification outputs are
recorded in \texttt{docs/03-windows-endpoint.md} in the project repository.

\subsection{Linux Endpoint}
\label{sec:impl-linux}

The Linux endpoint was provisioned by performing a clean Ubuntu 26.04 LTS
Desktop install from the official ISO. Unlike Wazuh and Windows, both of which
were imported from vendor-provided appliances, the Linux install was carried
out manually. This was a deliberate choice: the Linux install exposed
first-hand the number of small decisions (partitioning, timezone, initial
user creation) that shape a system's behaviour and that an appliance import
hides.

The Wazuh agent was installed on the Ubuntu VM using the manager's
wizard-generated \texttt{wget} and \texttt{dpkg} command, then enabled and
started as a systemd unit. It appeared as \texttt{Active} in the dashboard
with \texttt{agent.id} 002.

The critical configuration step on Linux was adding \textbf{auditd
command-execution monitoring}. Wazuh's default Linux agent configuration
reads system logs from \texttt{/var/log/syslog}, \texttt{/var/log/auth.log},
and, if present, \texttt{/var/log/audit/audit.log}. However, Ubuntu Desktop
does not include the Linux Audit framework by default, so \texttt{auditd} was
installed via \texttt{apt}, and audit rules were added to track every
invocation of the \texttt{execve} system call --- which is called whenever
any command is executed on the system:

\begin{quote}\small
\texttt{-a exit,always -F arch=b64 -S execve -k audit-wazuh-c}\\
\texttt{-a exit,always -F arch=b32 -S execve -k audit-wazuh-c}
\end{quote}

The \texttt{-k audit-wazuh-c} key identifies these events as
command-execution audits, so they can be selected in Wazuh queries. After
the rules were loaded with \texttt{augenrules --load}, the Wazuh agent was
restarted and began forwarding \texttt{audit.log} entries alongside the
default log sources.

At this point, both endpoints were producing detection-grade telemetry, and
the milestone identified in the project overview as M4 (both endpoints
delivering enhanced telemetry to the SIEM) was met. Full commands, output,
and verification screenshots are recorded in
\texttt{docs/04-linux-endpoint.md}.

\subsection{Ingestion Verification and Baselining}
\label{sec:impl-verify}

Once both endpoints were reporting, three verification checks were performed
before any attack simulation was attempted. First, the number of events per
minute in the dashboard's \textit{Threat Hunting} view was observed for each
endpoint over a five-minute window and cross-referenced against known
activity on that endpoint (for example, opening a text editor on the Linux
endpoint produced a visible spike in \texttt{execve} events). Second, the
alert-level distribution was checked: a healthy Wazuh deployment shows a
long tail of Level~1 to Level~3 informational events, a small number of Level~5
to Level~8 alerts, and rare Level~12+ high-severity events, which was
observed. Third, a manual query for \texttt{rule.description: "Sysmon"} on
the Windows agent, and \texttt{rule.description: "execve"} on the Linux
agent, was used to confirm that both endpoints' enhanced telemetry was not
merely being ingested but was actually reachable through the dashboard's
query interface.

With these verifications passed, the lab was considered operational and
ready for the attack simulations described in Section~\ref{sec:evaluation}.

\section{Evaluation}
\label{sec:evaluation}

\textit{This section is to be written. It will evaluate the lab against the
objectives set in Section~\ref{sec:aim} --- confirming that endpoints report
to the SIEM, that simulated security events generate the expected alerts, and
that each alert can be triaged --- supported by evidence and metrics drawn
from the running environment.}

\section{Discussion}
\label{sec:discussion}

\subsection{Overview}

This section revisits the three research questions posed in
Section~\ref{sec:rq}, reflects on the challenges encountered during the
build, and discusses the limitations of the lab in the context of a
professional SOC environment. The intent is to be honest about what the
project does and does not represent, so its value as a portfolio artefact
rests on defensible claims rather than on overstated equivalence.

\subsection{RQ1: Can a Credible SOC Environment Be Built with Only Free and Open-Source Tools?}

The evidence from the build supports a qualified yes. Wazuh, Sysmon, and
\texttt{auditd}, deployed on VirtualBox with two commodity endpoint
operating systems, produced a working SIEM environment capable of ingesting
detection-grade telemetry, applying a mature rule set out of the box, and
raising meaningful alerts within seconds of simulated attacker activity.
Every tool used is deployed in production by real organisations, and the
concepts practised in the lab --- agent-based collection, log parsing,
rule-based correlation, alert triage --- transfer directly to commercial
platforms.

The qualification is that "credible" here means credible for the purposes of
individual learning and portfolio demonstration, not for the purposes of
running a live security programme for a real organisation. Multi-tenant
scalability, high availability, long-term compliance retention, integration
with commercial threat intelligence feeds, and 24-hour analyst rotation are
all absent. What the lab does credibly represent is the daily operational
core of a SOC analyst's work: interpreting alerts, correlating events across
sources, classifying activity, and producing structured triage
documentation.

\subsection{RQ2: What Additional Configuration is Required Beyond Defaults?}

This turned out to be the most instructive question of the three. A default
installation of Wazuh, a default installation of Windows, and a default
installation of Ubuntu, connected together, produced far less useful
telemetry than a naive reading of the vendor documentation would suggest.
Three specific configuration additions were required to reach
detection-grade visibility:

\begin{enumerate}[label=\textbf{\arabic*.},leftmargin=*]
    \item On Windows, Sysmon had to be installed and its dedicated event
          channel (\texttt{Microsoft-Windows-Sysmon/Operational}) had to be
          explicitly forwarded by the Wazuh agent through an added
          \texttt{<localfile>} block in \texttt{ossec.conf}. Without this,
          the agent received only standard Windows Security event log
          entries, which do not include the process, network, and
          file-creation records that modern behavioural detection depends
          on.
    \item On Linux, \texttt{auditd} was not present in the default Ubuntu
          Desktop install. It had to be installed, and audit rules for the
          \texttt{execve} system call had to be added and loaded before the
          Wazuh agent could see command-execution events.
    \item Simulation 4 identified a third gap: the default Wazuh syscheck
          File Integrity Monitoring configuration on Windows runs on a
          twelve-hour scan cycle rather than in real-time. A modification
          to a security-sensitive file such as the hosts file will
          eventually be detected, but not with the latency a live SOC would
          need. Closing this gap requires adding
          \texttt{realtime="yes"} to the relevant
          \texttt{<directories>} entries.
\end{enumerate}

The recurring theme is that vendor defaults optimise for broad coverage with
low overhead, not for security-critical corner cases. A deployment that
relies only on defaults will have predictable blind spots, and the work of
detection engineering is largely the work of identifying and deliberately
closing those blind spots. This directly motivates the planned Project~2 on
detection engineering with MITRE ATT\&CK.

\subsection{RQ3: To What Extent Does This Lab Represent a Professional SOC?}

Realistically, the lab represents perhaps twenty per cent of a professional
SOC's operational surface, and does so along one axis: the individual
analyst's interaction with a SIEM. The remaining eighty per cent, honestly
enumerated, includes:

\begin{itemize}[leftmargin=*]
    \item \textbf{Scale.} A production SIEM ingests events at rates several
          orders of magnitude beyond what two lab endpoints produce. Alert
          fatigue, triage prioritisation, and tuning at scale are among
          the hardest and most important skills in a real SOC, and the lab
          cannot exercise them.
    \item \textbf{Team workflow.} There is no Tier~2 or Tier~3 escalation,
          no shift handover, no on-call rotation, no incident coordination
          across multiple analysts. The lab is a single-analyst
          environment.
    \item \textbf{Network detection.} All monitoring in the lab is
          host-based. A production SOC also relies on network-tier
          detection through IDS/IPS platforms such as Suricata or Zeek and
          on full packet capture. This is out of scope here and identified
          as future work in Section~\ref{sec:conclusions}.
    \item \textbf{Cloud and SaaS telemetry.} Modern SOCs increasingly work
          with cloud audit logs, SaaS activity, and identity provider
          events. The lab is entirely on-premises virtualised
          infrastructure.
    \item \textbf{Threat intelligence enrichment.} Alerts in the lab are
          not enriched with external context such as IP reputation,
          malware family attribution, or campaign intelligence. In a
          production environment this enrichment is a major driver of
          triage decisions.
    \item \textbf{Response tooling.} There is no EDR quarantine, no SOAR
          playbook automation, and no case management ticketing. Response
          actions in the lab were performed by hand.
\end{itemize}

Within its scope, however, the lab does honestly reflect one specific
piece of SOC work: what an analyst does in the first sixty seconds after
an alert appears. That is a small share of a SOC's total surface but a
significant share of a Tier~1 analyst's daily work, and it is the piece
that is hardest to practise without a working SIEM to practise on.

\subsection{Challenges Encountered}

Three challenges shaped the build in ways worth documenting for future
readers.

\textbf{Resource constraint on the host machine.} The 16~GB of RAM on the
host was tight for running three virtual machines concurrently under any
sustained load. In practice the two endpoints were alternated for heavy
testing while the Wazuh manager remained continuously running. This is a
legitimate constraint for a home lab of this class, but it would prevent
any experiment that required both endpoints active simultaneously (for
example, lateral movement simulations between them).

\textbf{Vulnerability-detector disk consumption.} During extended idle
periods between build sessions, the Wazuh manager's vulnerability-detector
cache expanded until it consumed the entire 25~GB manager disk. The
symptom was manager service failure with a critical configuration error
that traced back to \texttt{touch: No space left on device}. The fix was
to clear the \texttt{/var/ossec/queue/vd\_updater/} and
\texttt{/var/ossec/queue/vd/} directories, restart the indexer and
manager, and (in the longer term) either grow the manager disk or disable
the vulnerability-detector module in \texttt{ossec.conf}. This incident is
worth recording because it is the kind of operational maintenance a lab
owner learns only by encountering, and because it illustrates how a SIEM
component can fail silently between active use.

\textbf{Configuration file structure across versions.} Some steps in
public Wazuh guides refer to configuration paths or agent commands that
have changed between major versions. Working from the vendor's own
documentation for the specific version installed
(\texttt{4.14.5}) rather than from third-party tutorials generally
avoided this issue, and this experience reinforces a lesson from the
methodology: recording exact software versions at build time is not
optional.

\subsection{Honest Reflection on Portfolio Value}

The value of this project as a portfolio artefact is not that it makes the
author qualified to operate a production SOC on its own; it does not.
The value is that it provides evidence of three specific things which are
frequently claimed by graduate applicants but rarely demonstrated:

\begin{itemize}[leftmargin=*]
    \item \textbf{Hands-on SIEM operation}, rather than only textbook
          familiarity, evidenced by the deployed environment and the
          documented triage of real alerts.
    \item \textbf{Awareness of the gap between defaults and
          detection-grade configuration}, evidenced by the deliberate
          endpoint telemetry work and by the negative-result FIM
          simulation.
    \item \textbf{Structured technical documentation and analyst-style
          reporting}, evidenced by the incident reports and by this
          report itself, both of which are held to a standard that is
          suitable for handover to a colleague.
\end{itemize}

These are the specific competencies most difficult for a graduate applicant
to demonstrate from coursework alone, and they are the specific
competencies this project was designed to develop.

\section{Conclusions and Future Work}
\label{sec:conclusions}

\subsection{Summary}

This project designed, built, documented, and evaluated a Security
Operations Centre home lab using entirely free and open-source tooling.
The lab consists of a Wazuh SIEM manager and two monitored endpoints ---
one Windows and one Linux --- on an isolated virtualised network. Both
endpoints run deliberate detection-grade telemetry: Sysmon with the
SwiftOnSecurity baseline on Windows, and the Linux Audit framework with
\texttt{execve} monitoring on Linux. Four attack simulations covering
persistence, privilege escalation, credential access, and detection
configuration were executed and triaged against the built environment,
each producing either a documented true-positive alert cluster or a
deliberately documented negative result. Full analyst-style incident
reports were written for every simulation, and the entire build was
recorded as a reproducible step-by-step guide alongside a formal report
(this document).

The seven project objectives set out in Section~\ref{sec:aim} were all
met, and the four success criteria stated in the Initial Project Overview
were all satisfied.

\subsection{Answers to the Research Questions}

\textbf{RQ1.} A credible SOC monitoring environment can be built with only
free and open-source tools, provided that "credible" is scoped to
individual learning and portfolio demonstration rather than to running a
production security programme. The Wazuh stack in particular is genuinely
representative of a live SIEM in both operation and concept.

\textbf{RQ2.} Meaningful detection requires deliberate configuration
beyond defaults on both endpoint platforms: Sysmon plus a Wazuh-side
\texttt{<localfile>} block for Windows, and \texttt{auditd} with
\texttt{execve} rules for Linux. In addition, real-time File Integrity
Monitoring on security-critical paths requires an explicit
\texttt{realtime="yes"} directive, without which FIM operates on a
twelve-hour scan cycle. Vendor defaults optimise for coverage with low
overhead, not for security-critical corner cases, and closing those
corner cases is a large share of what detection engineering is.

\textbf{RQ3.} A single-machine virtualised lab represents perhaps twenty
per cent of a professional SOC's operational surface, and does so along
the axis most relevant to the Tier~1 analyst role: the first sixty
seconds after an alert appears. Scale, team workflow, network detection,
cloud telemetry, threat intelligence enrichment, and response automation
are all absent, and this is documented openly in the Discussion rather
than glossed over.

\subsection{Contribution}

The specific contribution of this project is a fully documented,
reproducible SOC home lab build using open-source tools, paired with a
set of triaged attack simulations that include both positive and
deliberately negative results, and with analyst-style incident reports
for each simulation. The project is deliberately positioned as the first
stage of a two-part progression: it establishes the monitoring foundation
on which a planned second project, focused on detection engineering with
MITRE ATT\&CK and Atomic Red Team, will build.

\subsection{Future Work}

Several extensions were identified during the project and are noted here
for future development. They fall into three groups.

\textbf{Detection engineering (Project~2).} The planned follow-on project
will focus on custom rule development, mapping detection coverage against
MITRE ATT\&CK, and using Atomic Red Team to execute a broader and more
systematic set of attack techniques. The two specific detection gaps
identified in this project --- the absence of an aggregate rule for
repeated failed authentication (Simulation~3) and the lack of real-time
FIM on the Windows hosts file (Simulation~4) --- are both natural starting
points for that work.

\textbf{Additional telemetry sources.} The current lab is host-based only.
Adding network-tier detection through Suricata or Zeek integrated with
Wazuh would extend coverage to lateral movement and command-and-control
traffic that host telemetry alone cannot see. Enabling the Wazuh
vulnerability-detector on a schedule (with the disk consumption issue
discussed in Section~\ref{sec:discussion} mitigated) would add
vulnerability visibility to the existing detection surface. Integrating a
threat intelligence feed for IOC enrichment would move the lab closer to
the enrichment workflows real SOCs depend on.

\textbf{Response automation and case management.} All response actions in
the current lab were performed by hand. Adding Wazuh's active-response
module, or integrating a lightweight SOAR platform such as Shuffle for
playbook automation, would demonstrate the response side of the SOC
lifecycle in the same practical way that this project demonstrates the
detection side. Adding a lightweight case management or ticketing layer
would let the analyst workflow be exercised end to end.

\subsection{Closing Note}

The core aim of this project was to close, for the author personally, the
gap between theoretical knowledge of SIEM operations and practical
hands-on experience of it. That gap has been closed, and the resulting
build and documentation are offered as both a self-contained portfolio
artefact and as a foundation for further work in detection engineering.
The environment stands ready for that follow-on work; the disciplines it
required to build --- deliberate telemetry design, structured triage,
honest documentation, and the willingness to run and record negative
tests --- are the same disciplines the follow-on work will demand at
greater depth.


% ----------------------------------------------------------------------
% References
% ----------------------------------------------------------------------
\newpage
\bibliographystyle{plain}
\bibliography{references}
\addcontentsline{toc}{section}{References}

\end{document}
